\scanpage{46}
\mspage{22}
Il faut cependant vérifier encore la condition MD sur la maquette $\underline G$ obtenue. Que $\underline G$ soit connexe résulte aussitôt de la connexité de $\underline G'$. Pour $\forall\alpha\in S$, il faut vérifier que
\[
2g_\alpha+\underbrace{\operatorname{card}I_\alpha+\operatorname{card}\widetilde A_\alpha}_{\widehat\nu_\alpha}\geq3
\]
\deleted{\unsure{or je dis qu'on aura mieux|ou je dis qu'on aura même}}

\deleted{ou}\inserted{i.e.} que $(g_\alpha,\widehat\nu_\alpha)$ est anabélien. Or
\[
\deleted{\operatorname{card}2g_\alpha+\operatorname{card}I_\alpha\geq3}
\]
\deleted{je dis}\inserted{on sait} que $(g_\alpha,\widehat\nu_\alpha)$ est le type numérique de la maquette $\underline G^{\deleted{\unsure{'}}}(\alpha)$, dont le graphe sous-jacent est la composante connexe de $G'_\bullet$ correspondant à $\alpha$, et l'ens. des ``pts marqués'' est $I'_\alpha\amalg\deleted{\unsure{\widetilde A}}\inserted{\widetilde B}'_\alpha$, \deleted{\unsure{et|à}} où $I'_\alpha=I'|S'_\alpha\simeq I_\alpha$, et $\widetilde B'_\alpha=\deleted{\unsure{\widetilde B|\widetilde A}}\inserted{\unsure{\widetilde A|\widetilde B}}'_\alpha\cap\widetilde B'$ est l'ens des flèches \unsure{de|$\in$} $\widetilde B'$ d'origine dans $S'_\alpha$, l'hom de projection
\[
\unsure{I_{(\alpha)}|I^{(\alpha)}}^{\deleted{\unsure{\ast|\alpha}}}\overset{\mathrm{def.}}{=}I'_\alpha\amalg\widetilde B'_\alpha\longrightarrow S'_\alpha
\]
étant induit par $I'\to S'$ et l'application origine. Donc, en vertu de la prop. du n° 3, on aura $(g_\alpha,\widehat\nu_\alpha)$ anabélien pourvu qu'on sache que $\underline G(\alpha)$ est un MD-graphe. Mais cette condition signifie que l'on a, pour $\forall\beta\in S'_{\deleted{\unsure{\beta|\alpha'}}\inserted{\alpha}}$
\mspage{45}
